\section{Deskripsi Persoalan}

Tugas besar ini membahas implementasi dan evaluasi dua keluarga arsitektur
\textit{deep learning}, yaitu \textit{Convolutional Neural Network} (CNN) dan
\textit{Recurrent Neural Network} (RNN), dalam dua persoalan berbeda. Bagian CNN
digunakan untuk persoalan klasifikasi citra, sedangkan bagian RNN dan LSTM
digunakan untuk persoalan \textit{image captioning}.
\subsection{Klasifikasi Citra dengan CNN}

Persoalan pertama adalah klasifikasi citra menggunakan dataset Intel Image
Classification. Dataset ini berisi gambar pemandangan alam dan lingkungan buatan
dengan enam kelas, yaitu \textit{buildings}, \textit{forest}, \textit{glacier},
\textit{mountain}, \textit{sea}, dan \textit{street}. Setiap gambar harus
diprediksi ke salah satu kelas tersebut berdasarkan pola visual yang muncul pada
citra.

CNN cocok untuk persoalan ini karena operasi konvolusi mampu menangkap pola
lokal seperti tepi, tekstur, bentuk, dan komposisi spasial. Pada tugas ini,
arsitektur CNN dilatih menggunakan Keras, kemudian bobot hasil pelatihan
digunakan kembali pada implementasi \textit{forward propagation} from \textit{scratch}.
Implementasi yang dibuat mencakup \code{Conv2D} dengan \textit{shared
parameters}, \code{LocallyConnected2D} dengan \textit{non-shared parameters},
pooling, global pooling, flatten, dense layer, serta fungsi aktivasi yang
diperlukan.

Eksperimen CNN dilakukan dengan beberapa variasi arsitektur, yaitu jumlah layer
konvolusi, jumlah filter, ukuran kernel, dan jenis pooling. Kinerja model
dievaluasi menggunakan \textit{macro F1-score} pada data uji agar performa setiap
kelas diperhitungkan secara seimbang. Selain itu, hasil \textit{forward
propagation} implementasi \textit{from scratch} dibandingkan dengan hasil Keras
untuk memastikan bahwa implementasi numeriknya sudah sesuai. Perbandingan antara
arsitektur \textit{shared} dan \textit{non-shared} juga dilakukan untuk melihat
pengaruh \textit{parameter sharing} terhadap akurasi, jumlah parameter, dan
efisiensi komputasi.

\subsection{Image Captioning dengan RNN dan LSTM}

Persoalan kedua adalah \textit{image captioning}, yaitu menghasilkan deskripsi
teks untuk sebuah gambar. Pada bagian ini digunakan dataset Flickr8k yang berisi
gambar beserta beberapa caption referensi untuk setiap gambar. Model harus
menghasilkan urutan kata yang menggambarkan isi gambar, sehingga persoalan ini
menggabungkan pemrosesan visual dan pemrosesan sekuensial.

Arsitektur yang digunakan mengikuti pendekatan \textit{encoder-decoder} dari
\textit{Show and Tell} oleh Vinyals et al. Gambar diproses oleh CNN pretrained
untuk menghasilkan \textit{feature vector}, lalu feature tersebut diproyeksikan
dan dimasukkan sebagai timestep awal decoder dengan skema \textit{pre-inject}.
Decoder berbasis Simple RNN atau LSTM kemudian menghasilkan caption kata demi
kata.

Pada tahap pelatihan, decoder dilatih menggunakan Keras dengan mekanisme
\textit{teacher forcing}. Input decoder terdiri atas feature vector gambar yang
telah diproyeksikan, token \code{<start>}, dan token-token caption sebelumnya,
sedangkan targetnya adalah token caption yang digeser satu posisi ke depan.
Bobot hasil pelatihan kemudian dapat dimuat ke implementasi \textit{from scratch}
yang terdiri atas embedding layer, Simple RNN cell, LSTM cell, dense projection,
dan dense output layer.

Eksperimen RNN dan LSTM bertujuan membandingkan pengaruh jumlah layer recurrent
dan ukuran hidden state terhadap kualitas caption. Evaluasi dilakukan
menggunakan metrik BLEU-4 dan METEOR, serta memperhatikan waktu eksekusi. Selain
itu, hasil decoder Keras dibandingkan dengan implementasi \textit{from scratch}
untuk memvalidasi kesesuaian \textit{forward propagation}. Perbandingan utama
antara Simple RNN dan LSTM digunakan untuk menganalisis perbedaan kemampuan
keduanya dalam mempertahankan informasi (memori) untuk jangka panjang pada urutan kata.
